\subsection{Provenance model}

The provenance chain is preserved for each \texttt{FinalRule} the RELATE stage emits, allowing a downstream reader to re-derive it from the source text. The chain has four links: 

\begin{itemize}
    \item \textbf{FinalRule $\to$ CanonicalRule.} Each \texttt{FinalRule} records the \texttt{canonical\_id} of the \texttt{CanonicalRule} it was scored from. The \texttt{CanonicalRule} stores the representative predicate; and is associated with its best-grounded member selected during CANONICALIZE, described in \pinktodo{Section 3.4.6, Canonicalize}. When needed, the verbatim text and the identifier of the paragraph it came from can be recovered on demand from the relational database.
    
    \item \textbf{CanonicalRule $\to$ NormalFinding.} Each \texttt{CanonicalRule} lists the \texttt{normal\_id}s of the \texttt{NormalFinding} records in its direction bin. Each \texttt{NormalFinding} records the MAP finding identifiers it was merged from, and its deduplicated source spans that connect the \texttt{NormalFinding} back to source text.
    
    \item \textbf{NormalFinding $\to$ source paragraph.} Each source span carries a citation identifier of the form  \texttt{$S_\textit{i}$|PMCID|te\_id}, where \texttt{te\_id} is the primary key of the text element the finding was derived from. The source texts lives in the relational database, and can be fetched on demand.
    
    \item \textbf{Source paragraph $\to$ source document.} Each paragraph row has a foreign key to its document row, which carries the paper's PMCID and document-level metadata.
\end{itemize}

\pinktodo{Write about the provenance chain.}

\pinktodo{Check if the empirical carry-rate of the pipeline is 100\%: we should empirically check if every final rule can be traced back to its source paragraph.}

The provenance chain supports downstream validation: any final rule can be opened, followed back to the \texttt{CanonicalRule} and the \texttt{NormalRule} records behind it, and the source paragraph, from which the findings were extracted. Then, the rule can be checked by a human reviewer against the source text, making the system auditable. However, provenance does not guarantee the correctness of the rule: a final rule might have a full provenance chain, and still encode an interpretation that a domain expert would reject. \pinktodo{Chapter 5, Discussion} discusses this distinction further. All in all, the final claim this thesis makes is this: the provenance chain establishes traceability, not clinical correctness.

